\section{Results II: Repeated Decoding and the \texttt{large-v3} Continuation Failure}

The most consequential unexpected result in the campaign arose when the
\texttt{large-v3} hard-speech benchmark was repeated.

The engine matrix first exposed the disagreement: one archived \texttt{large-v3} hard-split run
reported 7.69\% \wer{} with 184 insertions. We then ran a dedicated four-repeat probe over the
same 200 \texttt{test-other} utterances. Those controlled repeats produced 6.53\%, 6.07\%,
5.46\%, and 6.61\% \wer{}.

Error decomposition identifies what changed in that dedicated probe. Substitutions remained 87,
deletions remained 15, and correct-word hits remained 3,619 in every repeat, while insertions
ranged from 101 to 144. Thus the controlled repeated runs did not disagree about which spoken
words should be substituted or deleted. They disagreed about how much additional text to emit.
The earlier 7.69\% matrix run extends the campaign-wide observed insertion count to 184, but we
keep it separate from the controlled four-repeat range rather than silently pooling the two
experimental contexts.

This distinction matters for interactive dictation. An omitted or substituted word is typically
visible near the point at which it was spoken. A fluent inserted continuation can instead
introduce text that has no acoustic source and may extend beyond the intended utterance. Average
\wer{} treats all errors through the same aggregate numerator, but the user-facing failure is
qualitatively different.

\subsection{Decoder intervention}

Follow-up decoder experiments examined faster-whisper's temperature fallback and previous-text
conditioning. Simply forcing \texttt{temperature=0.0} did not provide a safe repair. In the
measured \texttt{large-v3} hard condition it could remove the fallback rescue path and produce
substantially worse output, including a 15.26\% \wer{} condition in the 2$\times$2 experiment.

Disabling \texttt{condition\_on\_previous\_text} produced a different result. In the measured
\texttt{large-v3} condition, it removed the observed runaway repetition and produced one
repeatable hypothesis set at 3.82\% \wer{} across the corresponding repetitions. However, the
paired corpus-average gain is not a statistically established general \wer{} improvement: the
estimated difference is approximately 1.05 \wer{} points with a 95\% interval spanning zero,
more utterances worsen than improve, and 95.7\% of the aggregate gain is concentrated in three
clips.

The appropriate interpretation is therefore \emph{tail-risk mitigation}, not average-accuracy
optimization.

The same intervention also behaves differently on smaller checkpoints. On
\texttt{test-other}, conditioning ON versus OFF produces 9.46\% versus 9.81\% for
\texttt{base.en}, 5.59\% versus 5.70\% for \texttt{small.en}, and byte-identical 5.51\%
output for \texttt{medium.en}. The direction therefore shrinks toward zero and then reverses at
\texttt{large-v3} in the measured ladder.

A later September follow-up also narrows an earlier interpretation of \texttt{tiny.en}. Five
complete decodes of the same 60-utterance \texttt{test-clean} subset produced byte-identical
hypotheses and 3.67\% \wer{} in all five runs. We therefore do not classify \texttt{tiny.en} as
generally corpus-level unstable. The strong repeated-corpus instability result in the present
evidence is specific to the measured \texttt{large-v3} condition.

These experiments illustrate why repeated decoding should be part of evaluation when a recognizer
contains fallback or sampling paths. A system may appear deterministic from a single benchmark
run while still exposing rare continuation behavior that is disproportionately costly in an
interactive setting.